% Môn: Vật lý
% Lớp: 11
% Chương: Chương I. Dao động
% Bài: L11C1 Bài 3. Vận tốc, gia tốc trong dao động điều hòa · Khai thác các thông số đặc trưng từ đồ thị li độ – thời gian
% ===== Câu 1 =====
% ID: L11C1B3-170-DS
% Mức: TH
\dangbt{Khai thác các thông số đặc trưng từ đồ thị li độ – thời gian}
\begin{ex}
% ID: L11C1B3-170-DS
% Mức: TH
Đồ thị li độ–thời gian có biên độ $4\,\text{cm}$, chu kì $1\,\text{s}$ và bắt đầu ở biên dương. Xác định tính đúng sai.
\choiceTF
{\True Biên độ bằng $4\,\text{cm}$.}
{\True Tần số bằng $1\,\text{Hz}$.}
{\True Tại $t=0$, vận tốc bằng 0.}
{\True Phương trình phù hợp là $x=4\cos(2\pi t/1)$ cm.}
\end{ex}

% ===== Câu 3 =====
% ID: L11C1B3-171-TLN
% Mức: TH
\begin{ex}
% ID: L11C1B3-171-TLN
% Mức: TH
Trên đồ thị li độ–thời gian, hai đỉnh dương liên tiếp ở $t=1$ s và $t=2$ s. Chu kì theo đơn vị s bằng bao nhiêu?
\shortans{1}
\end{ex}

% ===== Câu 6 =====
% ID: L11C1B3-172-TLN
% Mức: VD
\begin{ex}
% ID: L11C1B3-172-TLN
% Mức: VD
Trên đồ thị li độ–thời gian, hai đỉnh dương liên tiếp ở $t=1$ s và $t=3$ s. Chu kì theo đơn vị s bằng bao nhiêu?
\shortans{2}
\end{ex}

% ===== Câu 8 =====
% ID: L11C1B3-173-DS
% Mức: TH
\begin{ex}
% ID: L11C1B3-173-DS
% Mức: TH
Đồ thị li độ–thời gian cho biết $A=4$ cm, hai đỉnh dương liên tiếp cách nhau $1\,\text{s}$ và tại $t=0$ vật ở biên dương. Xác định $T,f,\omega$ và viết phương trình dao động.
\choiceTF

\end{ex}

% ===== Câu 9 =====
% ID: L11C1B3-174-TLN
% Mức: VD
\begin{ex}
% ID: L11C1B3-174-TLN
% Mức: VD
Trên đồ thị li độ–thời gian, hai đỉnh dương liên tiếp ở $t=1$ s và $t=5$ s. Chu kì theo đơn vị s bằng bao nhiêu?
\shortans{4}
\end{ex}

% ===== Câu 16 =====
% ID: L11C1B3-175-TN
% Mức: TH
\begin{ex}
% ID: L11C1B3-175-TN
% Mức: TH
Hai đỉnh dương liên tiếp trên đồ thị $x-t$ cách nhau $1\,\text{s}$. Chu kì bằng
\choice
{\True $1\,\text{s}$}
{$0.5\,\text{s}$}
{$2\,\text{s}$}
{$4\,\text{s}$}
\end{ex}
